\chapter{Marco teórico y estado del arte}
\label{cap:marco-teorico}

\section{Redes \emph{ad hoc} y multisalto}
\label{sec:ad-hoc}

Una red \emph{ad hoc} móvil (MANET) es una colección de nodos inalámbricos que se comunican sin infraestructura preexistente~\cite{perkins2003adhoc}. Cada nodo puede actuar como host y como router. En topologías multisalto, un paquete atraviesa nodos intermedios antes de alcanzar el destino final.

Los protocolos de enrutamiento MANET se clasifican en proactivos (mantienen tablas actualizadas continuamente, p.\ ej.\ OLSR) y reactivos (descubren rutas bajo demanda, p.\ ej.\ AODV, DSDV)~\cite{rfc3561}. En entornos con recursos limitados de batería y CPU ---como smartphones--- el equilibrio entre overhead de control y latencia de convergencia es crítico.

\section{Redes mesh en dispositivos móviles}
\label{sec:mesh-moviles}

Soluciones comerciales y de investigación han explorado mesh sobre Android mediante Wi-Fi Direct, Bluetooth o tunelización. Proyectos como Serval Mesh y Commotion Wireless demostraron viabilidad en escenarios comunitarios, pero suelen requerir \emph{root}, firmware modificado o arquitecturas distintas al perfil homogéneo GO+STA que propone PCT.

PCT se distingue por operar \textbf{exclusivamente en espacio de usuario}, sin \texttt{connect()} P2P para topología, y por separar explícitamente identidad lógica (UUID) de dirección IPv4 local.

\section{Wi-Fi Direct y SoftAP}
\label{sec:wifi-direct}

Wi-Fi Direct (Wi-Fi P2P) permite que dispositivos negocien un grupo con un Group Owner (GO) que actúa como SoftAP~\cite{wifidirect2010}. Android expone la API \texttt{WifiP2pManager} con operaciones \texttt{createGroup()}, \texttt{removeGroup()} y \texttt{requestGroupInfo()}.

La asociación STA legacy (\texttt{WifiNetworkSpecifier} desde API~29) permite que un dispositivo se conecte al SSID/PSK de otro GO como estación infraestructura, habilitando el patrón ``doble interfaz lógica'': GO propio para hijos downstream y STA upstream hacia el padre.

\section{DNS-SD y descubrimiento de servicios}
\label{sec:dns-sd}

DNS-Based Service Discovery (DNS-SD) sobre mDNS permite anunciar y descubrir servicios en redes locales~\cite{rfc6763}. Android Wi-Fi P2P integra DNS-SD mediante \texttt{WifiP2pDnsSdServiceInfo} y \texttt{addLocalService()}. PCT utiliza registros TXT compactos para transportar \texttt{node\_id}, rol, SSID y PSK del GO en la fase bootstrap (capa L1).

\section{Modelo OSI como referencia conceptual}
\label{sec:osi-emulado}

PCT no implementa una pila OSI completa; utiliza el modelo como \textbf{mapa de responsabilidades}:

\begin{table}[htbp]
  \centering
  \caption{Correspondencia capas OSI -- PCT}
  \label{tab:osi-pct}
  \begin{tabular}{@{}lll@{}}
    \toprule
    \textbf{Capa OSI} & \textbf{Capa PCT} & \textbf{Transporte} \\
    \midrule
    Física (L1)       & Radio / bootstrap & 802.11, DNS-SD \\
    Enlace (L2)       & Vecindad          & TCP :8765, :8766 \\
    Red (L3)          & Reenvío lógico    & DATA en :8766 \\
    Aplicación        & Mensajería usuario & Envelope \texttt{type=user} \\
    \bottomrule
  \end{tabular}
\end{table}

La regla normativa es que la capa de red \textbf{nunca} enruta por IP de destino final; la IPv4 (\texttt{next\_hop\_local\_ip}) es un dato de enlace para alcanzar al vecino inmediato.

\section{Estado del arte y posicionamiento de PCT}
\label{sec:estado-arte}

% TODO: ampliar con revisión bibliográfica sistemática (Scopus/IEEE)
Protocolos como B.A.T.M.A.N.~\cite{openwrt_batman} operan en capa~2 con visión global de vecinos por broadcast periódico. Babel~\cite{rfc8966} es un protocolo de enrutamiento vector-distancia para IPv6. Ambos asumen control sobre la interfaz de red a bajo nivel.

PCT adopta un enfoque de \textbf{capa de aplicación}: construye topología y enrutamiento sobre TCP confiable vecino-a-vecino, con identidad UUID y continuidad lógica (\texttt{epoch}, \texttt{session\_epoch}). Esta decisión sacrifica integración con el stack IP del sistema operativo a cambio de despliegue sin privilegios en Android comercial.

\section{Síntesis del capítulo}
\label{sec:sintesis-marco-teorico}

El marco teórico establece que las MANET multisalto en móviles requieren equilibrar overhead, identidad estable y restricciones de plataforma. Wi-Fi Direct GO + STA legacy + DNS-SD constituyen la base física de PCT; el modelo OSI emulado organiza las responsabilidades L1--L3 que se diseñan e implementan en capítulos posteriores.
